\par\medskip\noindent\textbf{Cómo leer la salida.}
{\small\setlength{\tabcolsep}{4pt}\renewcommand{\arraystretch}{1.15}
\begin{xltabular}{\linewidth}{@{}>{\raggedright\arraybackslash}p{0.24\linewidth}>{\raggedleft\arraybackslash}p{0.15\linewidth}>{\raggedright\arraybackslash}X@{}}
\toprule Indicador & Valor & Interpretación \\ \midrule \endfirsthead
\toprule Indicador & Valor & Interpretación \\ \midrule \endhead
\bottomrule \endlastfoot
\texttt{a}, \texttt{b} & 0,0225; 0,9732 & Parámetros del DCC con las tres series; persistencia 0,9957. \\
LR frente a CCC & 788,73 (2 g.l.) & Rechazo abrumador de correlaciones constantes. \\
Calma: S\&P-Nasdaq & 0,933 & Correlación promedio entre enero de 2006 y agosto de 2008. \\
Crisis: S\&P-Nasdaq & 0,961 & Sube algo en la crisis. \\
Calma / crisis: S\&P-Brent & $-0{,}053$ / 0,161 & Pasa de negativa a positiva: en la crisis acciones y petróleo cayeron juntos. \\
Calma / crisis: Nasdaq-Brent & $-0{,}066$ / 0,141 & El mismo patrón. \\
2017--2018 & 0,910; 0,103; 0,035 & Tras la crisis las correlaciones con el petróleo se moderan, pero no vuelven a ser negativas. \\
Volatilidad anual en la crisis & S\&P 50,0\,\%; Nasdaq 49,9\,\%; Brent 71,2\,\% & Volatilidades GARCH promedio, unas tres veces las normales: la diversificación falla justo cuando más se necesita. \\
Menor autovalor de $\bm R_t$ & 0,0295 & Positivo en todos los días: las matrices de correlación estimadas son válidas (definidas positivas). \\
\end{xltabular}}
